\documentclass[12pt]{amsart}
\usepackage{amssymb,amscd}
\usepackage{euscript}

\textheight25.5cm
\textwidth17.2cm
\hoffset=-2.5cm
\voffset=-2.8cm
\addtolength{\textheight}{1cm}
\leftmargin 3cm
\rightmargin 3cm

\def\u{\bold{u}}
\def\R{\mathbb R}
\def\Q{\mathbb Q}
\def\N{\mathbb N}
\def\Z{\mathbb Z}
\def\dist{\text{dist}}
\def\e{\varepsilon}


\def\QED{\hfill\hfill{\square}}  

\pagestyle{empty}

\begin{document}
%\setlength{\parindent}{0pt}


 {\bf Math  312H} \hskip.4cm  Lecture 13 \hskip.3cm 2/21/2025 \hskip.5cm 
 \vskip.3cm
 
 {\bf Bolzano-Weierstrass Theorem.\, Limits of subsequences.}   
 \vskip.25cm


       \vskip.35cm
       
  {\bf Last time:}\, If a sequence $(s_n)$ has a divergent subsequence, or two subsequences 
   \vskip.05cm 
\hskip2.4cm  converging to  different limits,\, then $(s_n)$ diverges.
\vskip.35cm

       
   {\bf Example.}\, Let $\,s_n=\cos\frac{n\pi}{100}$.\,    Then
   \vskip.1cm
\hskip2.2cm  $s_{200k}=\cos (2\pi k)=1$\, for all $k$,\, and so the subsequence $(s_{200k})$ converges to 1, and 
   \vskip.1cm
\hskip2.2cm  $s_{200k+100}=\cos (2\pi k+\pi)=-1$\, for all $k$,\, and so  $(s_{200k+100})$ converges to $-1$.
   \vskip.1cm
\hskip2.2cm  Since $(s_n)$ has two subsequences converging to different limits,\, it diverges.


\vskip.7cm

{\bf T/F}  \, Every sequence in $\R$ has a convergent subsequence.
   \vskip.1cm
\hskip1.2cm {\it False.}  \, Let $s_n=n$.\, Then every subsequence of $(s_n)$ diverges to $+\infty$.

 \vskip.4cm

{\bf Bolzano-Weierstrass Theorem.} \,
\vskip.08cm
Every  \underline{bounded} sequence in $\R$ has a convergent subsequence.

\vskip.3cm 

{\it Proof}. Let $(s_n)$ be a bounded sequence in $\R$. 
\vskip.06cm \hskip.6cm 
We will construct a Cauchy subsequence of $(s_n)$ inductively.\, Since $(s_n)$ is bounded, 
\vskip.06cm  \hskip.6cm 
there exists $M>0$\, such that for all $n\in \N$, \, $\,|s_n|\le M$\, and hence $\,-M\le s_n\le M$.
\vskip.06cm  \hskip.6cm 
Let $I_1=[-M, M]\,$ and let $\,n_1=1$.
\vskip.06cm  \hskip.6cm 
We divide $I_1$ into halves\, and denote a half with infinitely many terms of $(s_n)$ by $I_2$.
\vskip.06cm  \hskip.6cm 
We choose $n_2>n_1\,$ such that $\,s_{n_2}\in I_2$.
\vskip.06cm  \hskip.6cm 
Suppose that intervals $\,I_1\supset I_2\supset \dots \supset I_k\,$ and numbers $\,n_1<n_2<\dots <n_k$ are chosen.
\vskip.06cm  \hskip.6cm 
We divide $I_k$ into halves, denote a half with infinitely many terms of $(s_n)$ by $I_{k+1}$, 
\vskip.06cm  \hskip.6cm 
and choose $\,n_{k+1}>n_k\,$ such that $\,s_{n_{k+1}}\in I_{k+1}$.
\vskip.06cm  \hskip.6cm 
Thus we obtain a subsequence $(s_{n_k})$ of $(s_n)$.
\vskip.06cm  \hskip.6cm 
For each $N\in \N$,\,\,  $s_{n_k}\in I_N$\, for all $k\ge N$,\,\,  and $\,\, |I_N|=$ {\large $\frac{2M}{2^{N-1}}$}  $ \to 0$ as $N\to\infty$.
\vskip.06cm  \hskip.6cm 
It follows that $(s_{n_k})$ is Cauchy {\it (Explain)},\, and hence  $(s_{n_k})$ converges.  $\QED$

 \vskip.4cm
 
 {\bf Remarks}  
 \vskip.15cm 
$\,\,\bullet$ \,Using the same approach, one can show that any bounded sequence $(s_n)$ in $\R^2$ 
 \vskip.03cm
 \hskip.6cm  has a convergent subsequence. We observe that $(s_n)$ is contained in a square, 
  \vskip.03cm
  \hskip.6cm  and then at each step we divide the square into 4 equal squares.
 \vskip.15cm 
$\,\,\bullet$ \,In the book a different proof of Bolzano-Weierstrass Theorem is given.
 \vskip.03cm
 \hskip.6cm First, it is proved that any sequence in $\R$ has a monotone subsequence (Theorem 11.4),
 \vskip.03cm
 \hskip.6cm and then the theorem that every bounded monotone sequence in $\R$ converges is applied.

 \vskip.8cm

{\bf Limits of subsequences.}

\vskip.3cm

Let $(s_n)$ be a sequence in $\R$. A {\it subsequential limit}\, of $(s_n)$ is the limit of  a subsequence of $(s_n)$.
 \vskip.03cm
It can be a real number, $\infty$, or $-\infty$.
\vskip.4cm

   {\bf Examples.} \, $s_n=\cos\frac{n\pi}{4}+\frac1n.$ \,\, Subsequential limits:  $-1$, $-\frac{\sqrt{2}}{2}$,\, 0, $-\frac{\sqrt{2}}{2}$,\, 1.
    \vskip.25cm
 \hskip2.4cm  \, $s_n=(-n)^n.$ \,\, Subsequential limits: $-\infty$ and $\infty$.
 
 

\newpage 
   {\bf Theorem.}  Let $(s_n)$ be a bounded sequence in $\R$.
    \vskip.05cm
   Then $\,\limsup s_n$\, is the largest  subsequential limit of $(s_n)$,\,  and $\,\liminf s_n$\, is the smallest.
   
   \vskip.3cm 

{\it Proof for $limsup.$} \, Since $(s_n)$ is bounded, $\,b=\limsup s_n\in \R$.

   \vskip.2cm  
   
(1)  \,\,For each $\e>0$\, there exists $N\in \N$ such that\, $s_n<b+\e\,$ for all $n>N$.   

    \vskip.15cm 
    \hskip.8cm {\it Proof.}\,   Let $\e>0$.\, Since $\,b=\underset{N\to\infty}\lim \sup \{s_n:\, n>N\}$,\, there exists $N\in \N$\, such that
    
       \vskip.05cm 
    \hskip.8cm  $\sup \{s_n:\, n>N\}< b+\e$,\,\, and hence $\,s_n<b+\e\,$ for all $n>N$.
   

   \vskip.3cm  
   
(2) \,\,For each $\e>0$\, there are infinitely many $n$\, such that\, $s_n>b-\e$. 
      \vskip.15cm 
   \hskip.8cm {\it Proof.}\, Suppose that for some $\e>0$\, there are only finitely many such $n$.
\vskip.05cm 
   \hskip.8cm  Then there exists $N_1$\, such that $\,s_n\le b-\e\,$ for all $n> N_1$.
 Then for all $N>N_1$,\,\, 
     \vskip.05cm 
   \hskip.8cm  $v_N=\sup \{s_n:\, n>N\}\le b-\e$, \,and hence $\,b=\lim v_n\le b-\e$,\, which is impossible.
 
    \vskip.3cm  

(3)\,\,  There is a subsequence $(s_{n_k})$  of $(s_n)$ converging to $b$.
   
         \vskip.15cm 
   \hskip.8cm {\it Proof.}\, We inductively construct $(s_{n_k})$ such that for each $k\in \N$, \, $|s_{n_k}-b|<\frac1k$.
     \vskip.05cm 
   \hskip.8cm 
   Then it follows that $s_{n_k}\to b$. \, {\it (Explain.)} \,  By (1) and (2), or each $\e>0$\, 
     \vskip.05cm 
   \hskip.8cm  there are  infinitely many $n$\,  such that\, $b-\e< s_n<b+\e$,\, that is,  $\,|s_n-b|<\e$. 
   \vskip.05cm 
   \hskip.8cm   We choose $n_1$\, such that $\,|s_{n_1}-b|<1$.
   \vskip.05cm 
   \hskip.8cm Suppose that $\,n_1<\dots <n_k$\, are chosen so that $\,|s_{n_i}-b|<\frac{1}{i}$\, for $i=1, \dots, k.$
   
      \vskip.05cm 
   \hskip.8cm Since there are infinitely many $n$ such that $\,|s_{n}-b|<\frac{1}{k+1}$,\, we can choose $\,n_{k+1}>n_k\,$
      
      \vskip.05cm 
   \hskip.8cm such that $\,|s_{n_{k+1}}-b|<\frac{1}{k+1}$.
   
       \vskip.3cm  

(4)\,\, $(s_n)$ does not have a subsequence with $(s_{n_k})$\, with $\,\lim s_{n_k}>b$.

         \vskip.1cm 
   \hskip.8cm {\it Proof.}\, Exercise.  
   
\vskip.2cm  

Now the result follows from (3) and (4). $\QED$

   \vskip.6cm 
   
   {\bf Question.} How many subsequential limits can a sequence $(s_n)$ in $\R$  have?
\vskip.2cm 
  \hskip.8cm  For any $N\in \N$\, there is a sequence with exactly $N$ subsequential limits,\, 
  
   \hskip.8cm   for example\, $(s_n)=(1, 2, \dots, N, 1, 2, \dots, N,\dots).$
  \vskip.2cm   
   \hskip.8cm  Can  $(s_n)$ have infinitely many subsequential limits?
     \vskip.2cm   
{\bf Claim.}  There is a sequence in $\R$ such that \underline{every} $s\in \R$\, is its subsequential limit.
            \vskip.15cm 
   \hskip.8cm {\it Proof.}\, We construct a sequence $(s_n)$ that contains every rational number as follows.
         \vskip.05cm 
     \hskip.8cm First we list all fractions $\frac{m}{n}$, where $m,n\in \Z$ and $n\ne 0$,\, in an infinite table
           \vskip.05cm 
        \hskip.8cm and then we list the fractions following the diagonals
         \vskip.2cm 
 \hskip.8cm   0 $\rightarrow$ 1 $\rightarrow -1$ $\hskip.5cm 2 \hskip.4cm -2 \hskip.6cm 3 \hskip.4cm -3 \quad \dots$
 
      \hskip2.2cm  $\swarrow$   \hskip.5cm  $\swarrow$  \hskip.4cm  $\swarrow$  \hskip.5cm  $\swarrow$   \hskip.4cm  $\swarrow$
    
    \hskip1.8cm  $\frac12\hskip.4cm  -\frac12\hskip.7cm   \frac22\hskip.3cm -\frac22 \hskip.6cm \frac32 \hskip.4cm -\frac32 \quad \dots$
    
        \hskip2.2cm  $\swarrow$   \hskip.5cm  $\swarrow$  \hskip.4cm  $\swarrow$  \hskip.5cm  $\swarrow$   \hskip3.4cm So  $(s_n)=$
   
    \hskip1.8cm  $\frac13\hskip.6cm  -\frac13\hskip.7cm   \frac23\hskip.5cm -\frac23 \hskip.6cm \frac33 \hskip.55cm -\frac33 \quad \dots$
   \hskip1cm  $(0,\,1,\,-1,\, \frac12,\, 2,\, -\frac12,\, \frac13,\, -2,\, \frac22,\, -\frac13,\, \frac14, \, 3,\dots)$
\vskip.2cm
    
         \hskip1.8cm  . \, . \, . \, .\hskip5cm 
       
           \vskip.3cm  
           
          \hskip.8cm  For every $s\in \R$\, there is a subsequence $(s_{n_k})$ of $(s_n)$ converging to $s$.  
           
                      \vskip.05cm    
                 \hskip.8cm     This can be proved using an approach similar to (3) above (Homework). 
          
                      \vskip.05cm    
         \hskip.8cm   Since $(1,\,2,\,3,\, \dots)$ and $(-1,-2,-3,\, \dots)$ are subsequences of $(s_n)$, \, 
      $\infty$ and $-\infty$ are also
                 \vskip.05cm  
         \hskip.8cm  its  subsequential limits.  Thus the set of subsequential limits of $(s_n)$ is $\,\R\cup\{\infty, -\infty\}$.
    
\end{document}

%%%%%%%%%%%%%%%%%%%%%%%%%%
 For each $k\in \N$,\, there are infinitely many rational numbers in $(s-\frac1k, s+\frac1k)$, 
           \vskip.05cm    
         \hskip.8cm      
          and hence infinitely many $n$ such that $\,s_n\in (s-\frac1k, s+\frac1k)$.\,   As in (3), we can construct
      \vskip.05cm    
         \hskip.8cm  a subsequence $(s_{n_k})$ of $(s_n)$ such that for each $k\in \N$, \,
 $|s_{n_k}-b|<\frac1k$, \,  and hence $s_{n_k}\to s.$  \, 
